\section{Strategic dynamics: when everyone lives in a Bayesian world}
\label{sec:dynamics}

The chain now runs inside a population.  Theories determine committed
policies; policies determine which conditionals are visited; visited
conditionals determine dissonance and leverage; and refinements feed back
into theories.  This section characterizes where the process comes to rest.
The object of interest is not a strategic equilibrium but an epistemic
condition: a configuration of the population in which \emph{every agent
permanently inhabits a Bayesian world}---models fit, questions are dead, and
no represented calculation prices further inquiry above its cost.

\begin{lemma}[Finitely many refinements]
\label{lem:finite}
Each successful production strictly increases the number of cells of some
agent's partition, so along any path there are at most $n(|S|-1)$ successful
productions; almost surely there is a last one.
\end{lemma}

\begin{proof}
A strict refinement of a partition of a finite set adds at least one cell,
cell counts are bounded by $|S|$, and cumulative awareness
(\cref{ass:elective}) rules out coarsening.
\end{proof}

\begin{definition}[Self-confirming unawareness equilibrium]
\label{def:scue}
A profile $(\mathcal P_i,\mu_i,q_i,\sigma_i)_{i\in N}$ is a
\emph{self-confirming unawareness equilibrium} (SCUE) if, under the induced
play,
\begin{enumerate}[label=(\roman*)]
  \item each $\sigma_i$ is subjectively optimal
  (\cref{eq:subjective-optimality});
  \item each conjecture $q_i$ agrees with the play induced by $\sigma_{-i}$
  at every on-path cell;
  \item no dissonance occurs: each agent's active set is almost surely
  eventually nonempty, and each $\mu_i$ is supported on $\eta$-adequate
  theories; and
  \item no agent elects inquiry: $\lambda_iB_i<\gamma_i$, which holds in
  particular whenever $B_i=0$.
\end{enumerate}
\end{definition}

Conditions (i)--(ii) are subjective rationality with confirmed conjectures:
the self-confirming conditions of the fixed-language tradition
\citep{FudenbergLevine1993,KalaiLehrer1993b}.  Conditions (iii)--(iv) are the
new, epistemic ones: the language survives its own evidence, and the
question has no priced value.  A SCUE in which every partition is sufficient
for $\tau$ and beliefs concentrate near the truth is simply a
self-confirming equilibrium; the concept's interest lies entirely in the
other cases.

\begin{remark}[The grain of truth, restored endogenously]
\label{rem:grain}
Rational-learning convergence requires a grain of truth: the environment
absolutely continuous with respect to beliefs \citep{KalaiLehrer1993a}.
Under unawareness the condition can fail at the level of language---the
truth may be inexpressible---and statistical dissonance is precisely the
detector of that failure.  Each successful refinement restores the grain
\emph{on the path of play}: post-restoration, some expressible theory
matches the visited conditionals within tolerance.  A SCUE is a
configuration in which the grain of truth holds on-path for every agent and
no one has a priced reason to look for the off-path remainder.
\end{remark}

\begin{proposition}[Absorption]
\label{prop:absorption}
Suppose tests are consistent (\cref{lem:test-consistency}), tolerances
generic, conjectures empirically updated, and ties in
\cref{eq:subjective-optimality} broken by a fixed rule.  Then almost surely
there is a generation after which every agent's awareness is constant and
the population is in one of two absorbing regimes:
\begin{enumerate}[label=(\roman*)]
  \item a SCUE, in which every agent lives in a Bayesian world; or
  \item \emph{dissonant stasis} for a nonempty set of agents: dissonance
  recurs, but inquiry is never both elected and successful---because
  $\lambda_iB_i<\gamma_i$, or because elected production fails
  thereafter---and the affected agents carry committed policies forward with
  standing anomalies.
\end{enumerate}
If, whenever dissonance occurs with $B_i>0$, the election threshold is met
and $\phi_i$ succeeds with probability bounded away from zero, then regime
(ii) with $B_i>0$ is transient, and absorption into a SCUE or into
zero-leverage stasis is almost sure.  (Proof sketch in
\cref{app:absorption}.)
\end{proposition}

\begin{remark}[Rational anomaly tolerance]
\label{rem:anomaly}
Zero-leverage dissonant stasis is a recognizable condition, not a
pathology of the model: the agent knows its causal account is
broken---every expressible theory is rejected---yet the break disturbs no
represented action comparison, so repairing it is worth nothing.  Standing
anomalies tolerated by working scientists and managers have exactly this
structure, and the model locates them: anomaly tolerance is rational
precisely when pinned distinctions carry no robust action leverage.
\end{remark}

\begin{proposition}[On-path adequacy, and where error lives]
\label{prop:adequacy}
In a SCUE, every theory in the support of every $\mu_i$ is $\eta$-adequate
at every visited pair.  All surviving falsity is confined to what the
configuration does not expose: mechanism assignments at states pooled with
causally different states whose mixture imitates a single mechanism at the
visited profiles, and counterfactual claims at unvisited profiles.
\end{proposition}

\begin{proof}
Immediate from \cref{def:scue}(iii) and \cref{lem:test-consistency}; the
residual clause is the complement of the visited pairs' constraints.
\end{proof}

\begin{proposition}[Persistent performance differences]
\label{prop:persistence}
There exist primitives and absorbing configurations in which agents hold
different terminal partitions and theories, expected per-generation payoffs
differ strictly and permanently, and no disadvantaged agent has a priced
repair: for each such agent either $B_i=0$ or $\lambda_iB_i<\gamma_i$.
\end{proposition}

The proof is by the instances below, and the mechanism deserves emphasis.
Performance differences persist not because an equilibrium concept freezes
them but because the process that would erase them has died economically:
the disadvantaged configuration generates no dissonance, or generates no
leverage, or prices its question below the cost of asking.  Terminal
heterogeneity in awareness is path dependence in the strong sense---agents
with identical preferences and identical Bayesian competence, facing the
same objective system, end at different languages and different payoffs
because their histories posed different questions.

\begin{definition}[Diagnostic environments]
\label{def:diagnostic}
The primitives $(\Theta,\tau,\rho,u)$ are \emph{diagnostic} if, under every
profile sustainable in a SCUE, every expressible theory that is
payoff-relevantly false---false at some state where the truth would change
its holder's optimal action---is $\eta$-inadequate.
\end{definition}

\begin{proposition}[No payoff-relevant error in diagnostic environments]
\label{prop:diagnostic}
In a diagnostic environment, every SCUE is free of payoff-relevant error.
Payoff-relevant falsity survives only in environments that admit concealing
profiles, and it requires them to be sustained: collective error needs the
absence of \emph{both} corrective forces, strategic variety inside agents'
categories and environmental punishment of uniform misservice.
\end{proposition}

\begin{proof}
By \cref{def:scue}(iii), theories surviving in a SCUE are not
$\eta$-inadequate; in a diagnostic environment every payoff-relevantly false
theory is.  The second claim restates \cref{def:exposure}: a false theory
survives only under concealing profiles, and \cref{prop:exposure} makes
concealment necessary as well as sufficient.
\end{proof}

\begin{excont}
Set $\beta=1$ and suppose both firms acquired the column refinement after
the discovery generation and became certain of the false column assignment.
Each then plays the column policy; the firms pool at every state; every
payoff is one-half; and one-half is exactly what the column theory predicts
under pooled play.  The configuration is a SCUE with payoff-relevant error:
policies are subjectively optimal, conjectures confirmed, no dissonance ever
occurs (the profile is concealing), and $B_i=0$ for both firms---so the
configuration survives even a free, infallible inquiry technology.  The
warrant gap of \cref{sec:pivot} is permanent: unbounded experience
accumulates while each firm remains exactly half-warranted, the missing half
being the reflective weight on the world their shared language cannot
express.  Now let $\beta<1$.  In the first generation containing an
off-diagonal state, the pooled firms receive $\beta/2$ where their theory
predicts one-half: the demand side reveals the label, dissonance returns,
leverage turns positive, and the reflective posterior becomes degenerate.
Payoff-relevant self-confirmation is a knife-edge property of perfectly
forgiving demand---the $\beta<1$ environments are diagnostic---and the
destroyed value $1-\beta$ at each misserved episode is itself the corrective
signal: what demand withholds from the firms, their epistemics receives as
evidence.  Finally, persistence: with $\beta=1$, let firm~1 hold the row
refinement and the true assignment while firm~2 remains coarse with
$\gamma_2>\lambda_2B_2$.  Firm~1 wins both off-diagonal states and splits
the rest, earning three-quarters to firm~2's one-quarter, forever; firm~2
sits in dissonant stasis, its question priced below its cost.  Every route
out---strategic variety, demand harshness, cheaper inquiry---is exactly one
of the model's named forces.
\end{excont}
